% from ega1-6-fr.tex
\oldpage[I]{142}
tout revient à démontrer la proposition lorsque $Y=D(f)$ avec $f\in A$. Mais
alors $Y$ est un schéma affine dont l'anneau est isomorphe à $A_f$
(\hyperref[I.1.3.6-fr]{1.3.6})~; comme $A$ est noethérien
(\hyperref[I.6.1.3-fr]{6.1.3}), il en est de même de $A_f$.
\end{proof}

\begin{env}[6.1.5]
\label{I.6.1.5-fr}
On notera que le \emph{produit} de deux $S$-préschémas noethériens n'est pas
nécessairement noethérien, même si ces préschémas sont affines, car le produit
tensoriel de deux algèbres noethériennes n'est pas nécessairement un anneau
noethérien (cf. (\hyperref[I.6.3.8-fr]{6.3.8})).
\end{env}

\begin{proposition}[6.1.6]
\label{I.6.1.6-fr}
Si $X$ est un préschéma noethérien, le Nilradical $\mathscr{N}_X$ de
$\mathscr{O}_X$ est nilpotent.
\end{proposition}

On peut en effet recouvrir $X$ par un nombre fini d'ouverts affines $U_i$ et il
suffit de prouver qu'il existe des entiers $n_i$ tels que
$(\mathscr{N}_X|U_i)^{n_i}=0$~; si $n$ est le plus grand des $n_i$, on aura
alors $\mathscr{N}_X^n=0$. On est donc ramené au cas où
$X=\operatorname{Spec}(A)$ est affine, $A$ étant un anneau noethérien~; en
vertu de (\hyperref[I.5.1.1-fr]{5.1.1}) et de
(\hyperref[I.1.3.13-fr]{1.3.13}), il suffit d'observer que le nilradical de
$A$ est nilpotent (\cite{I-11}, p.~127, cor.~4).

\begin{corollary}[6.1.7]
\label{I.6.1.7-fr}
Soit $X$ un préschéma noethérien~; pour que $X$ soit un schéma affine, il faut
et il suffit que $X_{\mathrm{red}}$ le soit.
\end{corollary}

Cela résulte de (\hyperref[I.6.1.6-fr]{6.1.6}) et
(\hyperref[I.5.1.10-fr]{5.1.10}).

\begin{lemma}[6.1.8]
\label{I.6.1.8-fr}
Soient $X$ un espace topologique, $x$ un point de $X$, $U$ un voisinage ouvert
de $x$ n'ayant qu'un nombre fini de composantes irréductibles. Alors il existe
un voisinage $V$ de $x$ tel que tout voisinage ouvert de $x$ contenu dans $V$
soit connexe.
\end{lemma}

Soient $U_i$ ($1\leq i\leq m$) les composantes irréductibles de $U$ ne
contenant pas $x$~; le complémentaire dans $U$ de la réunion des $U_i$ est un
voisinage ouvert $V$ de $x$ dans $U$, donc aussi dans $X$~; c'est d'ailleurs
le complémentaire dans $X$ de la réunion des composantes irréductibles de $X$
qui ne contiennent pas $x$ (\hyperref[0.2.1.6-fr]{0, 2.1.6}). Soit alors $W$
un voisinage ouvert de $x$ contenu dans $V$. Les composantes irréductibles de
$W$ sont les traces sur $W$ des composantes irréductibles de $U$ qui
rencontrent $W$ (\hyperref[0.2.1.6-fr]{0, 2.1.6}), donc ces composantes
contiennent $x$~; comme elles sont connexes, il en est de même de $W$.

\begin{corollary}[6.1.9]
\label{I.6.1.9-fr}
Un espace topologique localement noethérien est localement connexe (ce qui
entraîne entre autres que ses composantes connexes sont ouvertes).
\end{corollary}

\begin{proposition}[6.1.10]
\label{I.6.1.10-fr}
Soit $X$ un espace topologique localement noethérien. Les conditions suivantes
sont équivalentes~:
\begin{enumerate}
  \item[a)] Les composantes irréductibles de $X$ sont ouvertes.
  \item[b)] Les composantes irréductibles de $X$ sont identiques à ses
    composantes connexes.
  \item[c)] Les composantes connexes de $X$ sont irréductibles.
  \item[d)] Deux composantes irréductibles distinctes de $X$ ne se rencontrent
    pas.
\end{enumerate}

Enfin, si $X$ est un préschéma, ces conditions équivalent aussi à~:
\begin{enumerate}
  \item[e)] Pour tout $x\in X$, $\operatorname{Spec}(\mathscr{O}_x)$ est
    irréductible (autrement dit le nilradical de $\mathscr{O}_x$ est premier).
\end{enumerate}
\end{proposition}

Il est immédiat que a) entraîne b), car un espace irréductible est connexe, et
a) entraîne que les composantes irréductibles de $X$ sont des ensembles à la
fois ouverts et fermés. Il est trivial que b) entraîne c)~; inversement, un
ensemble fermé $F$ contenant
